\chapter{Záver}

V predkladanom projektovom reporte sme venovali pozornosť logistickej 
sieti leteckých spojení na celom svete. Analyzovaná sieť obsahovala $3214$ letísk
a dokopy $36 \> 907$ orientovaných hrán -- jednotlivých leteckých spojení.

Detailnejšia analýza celej siete preukázala relatívne nízku hustotu na úrovni $0.357 \%$, čo preukazuje
priamu dostupnosť letísk na významne nízkej úrovni. Na druhú stranu sa preukazuje vysoká miera tranzitivity siete ($25 \%$)
a taktiež existencia významne veľkého komponentu súvislosti (podiel $> 99 \%$).

Samotná sieť vykazuje, oproti vybraným náhodným sieťovým modelom, vysoký koeficient zhlukovania, čo naznačuje, i napriek nízkej miere 
priamych leteckých spojení organizovanú štruktúrovanosť siete. Berúc do úvahy tieto závery, je možné pozorovať nízku hodnotu priemerného
počtu minimálnych nutných prestupov ($2.98$), čo preukazuje, že daná sieť je reprezentantom siete \textbf{malého sveta}.

Sieť preukazuje hodnotu asortativity blízku $0$, čo naznačuje, že sieť nevykazuje výraznejšiu tendenciu spájať podobne vyťažené letiská -- veľké
hubové letiská sa rovnako bežne spájajú s menšími regionálnymi letiskami ako s inými hubmi. Okrem iného sieť je takmer neorientovaná, s mierou reciprocity necelých 
$98 \%$.

Sieť obsahovala dokopy $718$ listov -- vrcholov so stupňom $1$. Podobne bolo identifikovaných $333$ artikulácií, pričom sa primárne jedná o letiská v oblastiach napr. USA, Kanady či Južnej Ameriky.
V rovnakom duchu bolo identifikovaných $757$ mostových spojení, avšak v drvivej väčšine sa jednalo o spojenia, ktorého odstránenie viedlo k odstrihnutiu jedného letiska. 
Zaujímavým pozorovaním je to, že artikulačnými vrcholmi sú aj vysoko vyťažované medzinárodné letiská (napr. Tokyo alebo Denver), ktorých odstránením docháza k tvorbe aj viac ako $10$ nových komponentov súvislosti.
Tu sa preukazuje kritickosť veľkých letísk v kontexte súvislosti celej siete.

